% GENERATED FILE. Edit canonical lessons, then run npm run generate:books.
% canonical-source-hash: fnv1a64:280b87519268ce11
% canonical-lessons: GE-C288-zusammenstossen, GE-C288-zusammenstoss, GE-C288-prallen, GE-C288-verungluecken, GE-C288-schaden

\chapter{Ein Zusammenstoß}
\label{ch:ge-ein-zusammensto}
\hlchaptermodality{288}

\begin{chapteropening}
\textbf{By the end of this chapter:} \emph{I can report a road accident: to collide, a collision, to crash into, to have an accident and the damage.}

Complete the last lesson of chapter 288: I can report a road accident: to collide, a collision, to crash into, to have an accident and the damage.
\end{chapteropening}

\section[zusammenstoßen]{zusammenstoßen — to collide}

\label{lesson:GE-C288-zusammenstossen}



\begin{quote}
Before the new one: say the German for to disappear, then the German for to flee.
\end{quote}

\subsection*{You'll want to know: zusammenstoßen}
\textbf{zusammenstoßen} — \textquotedblleft{}to collide\textquotedblright{}.

\begin{quote}
Zwei Autos stießen an der Kreuzung zusammen.
\end{quote}

\begin{quote}
\emph{Two cars collided at the crossroads.}
\end{quote}

Separable, Perfekt with \emph{sein}: \emph{Sie sind zusammengestoßen}.

\begin{cousinweb}
\emph{zusammenstoßen} joins \emph{zusammen} \textquotedblleft{}together\textquotedblright{} and \emph{stoßen} \textquotedblleft{}to push, to bump\textquotedblright{}, from Old High German \emph{stōzan}, Proto-Germanic \emph{stautaną} (EtymDB).
\end{cousinweb}

\begin{grammarlens}[title={What you've built}]
One more word for telling, explaining and describing at B1.
\end{grammarlens}

\subsection*{Your turn}
\begin{itemize}
\raggedright
  \item \emph{Say it:} \emph{zusammenstoßen}
  \item \emph{Say it:} \emph{zusammenstoßen}, once more
  \item \emph{Say it:} a sentence of your own with \emph{zusammenstoßen}
  \item \emph{Recall:} say the German for to disappear, then the German for to flee, then say \emph{zusammenstoßen} again
\end{itemize}

\begin{tcolorbox}[breakable,colback=teal!4,colframe=teal!35!black,title={Before you move on}]
What does zusammenstoßen mean? (\textquotedblleft{}To collide\textquotedblright{}.) Say it once more.
\end{tcolorbox}
\section[der Zusammenstoß]{der Zusammenstoß — the collision, the crash}

\label{lesson:GE-C288-zusammenstoss}



\begin{quote}
Before the new one: say the German for to flee, then the German for to collide.
\end{quote}

\subsection*{You'll want to know: der Zusammenstoß}
\textbf{der Zusammenstoß} — \textquotedblleft{}the collision, the crash\textquotedblright{}.

\begin{quote}
Bei dem Zusammenstoß ging ein Fenster kaputt.
\end{quote}

\begin{quote}
\emph{A window broke in the collision.}
\end{quote}

Plural: \emph{die Zusammenstöße}.

\begin{cousinweb}
\emph{Zusammenstoß} is built on \emph{zusammenstoßen} \textquotedblleft{}to collide\textquotedblright{}, from \emph{stoßen} \textquotedblleft{}to push\textquotedblright{}, Proto-Germanic \emph{stautaną} (EtymDB).
\end{cousinweb}

\begin{grammarlens}[title={What you've built}]
One more word for telling, explaining and describing at B1.
\end{grammarlens}

\subsection*{Your turn}
\begin{itemize}
\raggedright
  \item \emph{Say it:} \emph{der Zusammenstoß}
  \item \emph{Say it:} \emph{der Zusammenstoß}, once more
  \item \emph{Say it:} a sentence of your own with \emph{der Zusammenstoß}
  \item \emph{Recall:} say the German for to flee, then the German for to collide, then say \emph{der Zusammenstoß} again
\end{itemize}

\begin{tcolorbox}[breakable,colback=teal!4,colframe=teal!35!black,title={Before you move on}]
What does der Zusammenstoß mean? (\textquotedblleft{}The collision, the crash\textquotedblright{}.) Say it once more.
\end{tcolorbox}
\section[prallen]{prallen — to crash (into), to hit hard}

\label{lesson:GE-C288-prallen}



\begin{quote}
Before the new one: say the German for to collide, then the German for the collision.
\end{quote}

\subsection*{You'll want to know: prallen}
\textbf{prallen} — \textquotedblleft{}to crash (into), to hit hard\textquotedblright{}.

\begin{quote}
Das Auto prallte gegen einen Baum.
\end{quote}

\begin{quote}
\emph{The car crashed into a tree.}
\end{quote}

Perfekt with \emph{sein}: \emph{Es ist gegen die Wand geprallt}.

\begin{cousinweb}
\emph{prallen} is usually traced to Middle High German \emph{prellen} \textquotedblleft{}to bounce off\textquotedblright{}; EtymDB lists no parent.
\end{cousinweb}

\begin{grammarlens}[title={What you've built}]
One more word for telling, explaining and describing at B1.
\end{grammarlens}

\subsection*{Your turn}
\begin{itemize}
\raggedright
  \item \emph{Say it:} \emph{prallen}
  \item \emph{Say it:} \emph{prallen}, once more
  \item \emph{Say it:} a sentence of your own with \emph{prallen}
  \item \emph{Recall:} say the German for to collide, then the German for the collision, then say \emph{prallen} again
\end{itemize}

\begin{tcolorbox}[breakable,colback=teal!4,colframe=teal!35!black,title={Before you move on}]
What does prallen mean? (\textquotedblleft{}To crash (into), to hit hard\textquotedblright{}.) Say it once more.
\end{tcolorbox}
\section[verunglücken]{verunglücken — to have an accident}

\label{lesson:GE-C288-verungluecken}



\begin{quote}
Before the new one: say the German for the collision, then the German for to crash.
\end{quote}

\subsection*{You'll want to know: verunglücken}
\textbf{verunglücken} — \textquotedblleft{}to have an accident\textquotedblright{}.

\begin{quote}
Mein Onkel verunglückte mit dem Fahrrad.
\end{quote}

\begin{quote}
\emph{My uncle had an accident on his bike.}
\end{quote}

Perfekt with \emph{sein}: \emph{Er ist verunglückt}. Built on \emph{das Unglück}, misfortune.

\begin{cousinweb}
\emph{verunglücken} is built on \emph{Unglück} \textquotedblleft{}misfortune\textquotedblright{}, Middle High German \emph{unglück} (EtymDB), from \emph{Glück} \textquotedblleft{}luck\textquotedblright{}, Middle High German \emph{gelücke}.
\end{cousinweb}

\begin{grammarlens}[title={What you've built}]
One more word for telling, explaining and describing at B1.
\end{grammarlens}

\subsection*{Your turn}
\begin{itemize}
\raggedright
  \item \emph{Say it:} \emph{verunglücken}
  \item \emph{Say it:} \emph{verunglücken}, once more
  \item \emph{Say it:} a sentence of your own with \emph{verunglücken}
  \item \emph{Recall:} say the German for the collision, then the German for to crash, then say \emph{verunglücken} again
\end{itemize}

\begin{tcolorbox}[breakable,colback=teal!4,colframe=teal!35!black,title={Before you move on}]
What does verunglücken mean? (\textquotedblleft{}To have an accident\textquotedblright{}.) Say it once more.
\end{tcolorbox}
\section[der Schaden]{der Schaden — the damage}

\label{lesson:GE-C288-schaden}



\begin{quote}
Before the new one: say the German for to crash, then the German for to have an accident.
\end{quote}

\subsection*{You'll want to know: der Schaden}
\textbf{der Schaden} — \textquotedblleft{}the damage\textquotedblright{}.

\begin{quote}
Der Schaden am Auto war groß.
\end{quote}

\begin{quote}
\emph{The damage to the car was great.}
\end{quote}

Plural: \emph{die Schäden}. \emph{Schade!} — what a pity — comes from the same word.

\begin{cousinweb}
\emph{Schaden} is usually traced to Old High German \emph{scado} \textquotedblleft{}damage\textquotedblright{}, which EtymDB derives from Proto-Germanic \emph{skaþô}; EtymDB lists English \emph{scathe} as a cousin.
\end{cousinweb}

\begin{grammarlens}[title={What you've built}]
One more word for telling, explaining and describing at B1.
\end{grammarlens}

\subsection*{Your turn}
\begin{itemize}
\raggedright
  \item \emph{Say it:} \emph{der Schaden}
  \item \emph{Say it:} \emph{der Schaden}, once more
  \item \emph{Say it:} a sentence of your own with \emph{der Schaden}
  \item \emph{Recall:} say the German for to crash, then the German for to have an accident, then say \emph{der Schaden} again
\end{itemize}

\begin{tcolorbox}[breakable,colback=teal!4,colframe=teal!35!black,title={Before you move on}]
What does der Schaden mean? (\textquotedblleft{}The damage\textquotedblright{}.) Say it once more.
\end{tcolorbox}
